% ============================================================
% 第1章：绪论
% ============================================================

\section{研究背景与问题定义}
电子元器件选型在硬件产品研发全生命周期中具有基础性和关键性作用。以 DC-DC 电源模块设计为例，德州仪器（Texas Instruments）、亚德诺半导体（Analog Devices）、意法半导体（STMicroelectronics）、微芯科技（Microchip Technology）等数十家制造商提供了数万种相关型号，工程师通常需要逐项核查并交叉比对其输入输出电压范围、最大输出电流、开关频率、工作温度等级、封装约束、库存水平和生命周期状态等多维参数；同时，还需查阅数据手册以验证外围元件的设计公式，并综合评估供应链稳定性 \cite{pecht2008parts}。该流程通常耗时数小时以上，且主要存在三方面瓶颈：第一，工程知识检索高度碎片化，设计规范分散于数据手册、应用笔记和工程师个人经验中；第二，供应链风险缺乏系统化的量化评估手段 \cite{chen2021supplyrisk}；第三，选型结论和输出文档格式不够规范，难以直接支撑工程评审或产品生命周期管理平台的集成应用。

针对上述问题，若能充分利用元器件管理平台的结构化物料数据与参考设计资源，并结合智能化推理技术，有望提升选型过程的自动化程度、风险识别能力和文档生成规范性。硬禾科技旗下的 eZ-PLM 是面向电子元器件全生命周期管理的开放平台（www.ezplm.cn/www.ezmpl.com），涵盖超过 110 万条在产物料数据与参考设计资源，并支持进行实时器件查询和设计案例检索，为器件检索、参考设计复用和供应链信息获取提供了数据基础。在此基础上，近年来大语言模型（LLM）与智能体（Agent）技术在工程决策辅助领域也展现出较强潜力 \cite{xi2023rise,li2023llmagent}，其自然语言理解、任务分解与交互推理能力可与领域规则引擎形成互补。因此，将 eZ-PLM 的结构化物料数据与 LLM/Agent 的智能推理能力相结合，有望将传统人工选型流程转化为智能化、可追溯的自动化决策过程。
% [插图提醒：图1-1 系统应用场景示意图——展示工程师从自然语言需求输入到获得BOM选型清单、风险评估报告、拓扑分析报告的端到端流程，左侧为用户输入框，中间为Pipeline五阶段，右侧为三类输出报告]

\section{现有方案的不足}

现有的元器件选型辅助方案大致可分为基于传统工业软件的方案和基于通用大语言模型的方案两类，但各自在工程实用性方面均存在显著不足。

\subsection{传统PLM/ERP方案的局限}

以ERP和PLM平台为代表的传统工业软件（如SAP Integrated Product Development、Oracle NetSuite、PTC Windchill等）是当前企业元器件选型管理的主流工具。这些平台通常提供基于表单的参数过滤、库存查询、型号管理和BOM导出等基础功能，但在智能化程度方面存在以下关键局限：

\textbf{（1）缺乏自然语言交互能力。}工程师必须将设计需求手动映射为平台特定的筛选字段和参数范围，无法直接以自然语言描述选型意图。例如，“需要一个12V转5V、3A的车规级降压方案”这样的需求，需拆解为“类别=DC-DC转换器”“拓扑=Buck”“输入电压≈12V”“输出电压=5V”“输出电流≥3A”“温度等级=车规”等字段逐一填入筛选表单。这一映射过程不仅耗时，还容易因人为疏漏导致约束条件不完整。

\textbf{（2）无应用场景适配度评估。}传统平台仅根据参数范围匹配返回符合条件的器件列表，不具备对器件的应用场景适配度进行智能推荐的能力。例如，两款均满足“输入电压4.5–40V、输出5V/3A”的器件，一款为面向通信设备的低噪声型号，另一款为面向工业电机驱动的宽温型号——传统平台无法区分两者在实际应用场景中的适配差异，工程师仍需凭经验做出最终判断。

\textbf{（3）供应链风险评估能力薄弱。}多数PLM平台仅能展示库存数量和生命周期状态等静态信息，缺乏对供应商集中度\cite{pecht2002electronic}、地缘政治风险、交期波动等多维度供应链风险的量化评估和分级预警能力\cite{goel2021supply}。器件停产（EOL/Obsolescence）的早期预警与主动替代推荐功能尤为缺失\cite{sandborn2013forecasting}。

\subsection{通用LLM方案的幻觉与不可溯源性}

以ChatGPT、DeepSeek等为代表的通用大语言模型在自然语言理解与工程知识表达方面展现了较强能力，但在元器件选型这一高可靠性场景中面临根本性挑战——幻觉（hallucination）问题\cite{wang2023survey}。具体表现为三种典型形式：

\textbf{（1）型号编造。}当被询问某款电源模块是否适配特定输入/输出条件时，模型可能编造一个看似合理但实际不存在的器件型号，并附上伪造的电压、电流和封装参数。例如，在没有检索到真实器件数据库的情况下，模型可能生成“TPS12345-Q1, 4.5V–36V输入, 5V/3A输出, AEC-Q100认证”这类完全虚构的推荐结果。

\textbf{（2）参数捏造。}模型可能在引用真实器件型号的同时，对该器件的关键参数进行错误描述——例如将LM2596S-5.0的输出电流误报为5A（实际为3A），或声称某器件支持AEC-Q100认证而实际并未取得该认证。这类错误对于依赖参数准确性的工程决策是致命的。

\textbf{（3）不可溯源。}通用LLM的输出无法追溯到具体的数据来源——无法说明某条推荐是基于哪家制造商的数据手册、哪份应用笔记或哪条API返回的实时库存记录。这导致工程师无法验证选型结论的可靠性，也无法满足ISO 9001等质量管理体系对设计决策可追溯性的要求\cite{iso31000}。

\subsection{两类方案之间的能力断层}

综上所述，当前元器件选型领域存在一个显著的能力断层：传统PLM/ERP方案具备真实的物料数据和结构化的参数信息，但缺乏智能化的需求理解与风险评估能力；通用LLM方案具备自然语言交互与知识表达能力，但缺乏真实数据支撑和结果可溯源性。能够将eZ-PLM平台的真实物料数据、参考设计资源与LLM的智能推理能力有机结合，并在工程规则引擎的约束下保障输出的严谨性与可追溯性，正是本系统设计的出发点\cite{lewis2020rag,gao2023retrieval}。


\section{本文贡献}

本文面向eZ-PLM平台，设计并实现了一套基于LLM与工程规则引擎协同推理的智能元器件选型与风险评估Agent系统。系统的主要技术贡献包括以下五个方面：

一是提出了五阶段线性流水线（Pipeline）架构，将"自然语言需求解析—实时API器件检索—混合智能评分—证据链构建—标准化报告生成"串联为端到端自动化流程，并辅以基于LangChain框架的推理行动（ReAct）Agent交互模式\cite{yao2022react}，同时支持批量处理与多轮对话两种使用场景；

二是设计了四级容错需求解析链，依次通过LLM语义解析、规则化类别覆盖、电压比较方向兜底和温度范围正则匹配来修正mA单位误识别、否定句式漏判、LLM拓扑幻觉等工程场景中的典型解析缺陷；

三是构建了以eZ-PLM参考设计数据为上下文约束的LLM混合评分机制，实现了"参数匹配＋供应链稳定性＋应用适配度＋设计成熟度"四维加权评分，同时支持无参考设计时自动降级为纯规则双维评分，保证了评分结果的可降级性与可比性；

四是建立了规则引擎事实验证与LLM工程叙述分离的双层风险评估架构——规则侧负责检测供应商集中度、零库存、停产等客观风险，LLM侧仅对已验证的风险条目进行工程背景阐述，确保了风险报告的严谨性与数据可溯源性；

五是定义了PartIR、ScoreBreakdown、TopologyIR、RiskIR等标准化中间表示（IR）数据模型，三类工程输出文档（物料清单（BOM）选型清单、供应链风险评估报告、电路拓扑分析报告）的格式可直接用于工程评审及eZ-PLM平台对接。
